\section{Introduction}

Dadas el conjunto de ecuaciones de Estructura Estelar.

\begin{equation*}
	\left.\begin{aligned}
		\diff{P}{r} & = -\frac{GM\rho}{r^{2}}.                \\
		\diff{M}{r} & = 4\pi r^{2}\rho.                       \\
		P           & = K{\left(\rho\right)}^{1+\frac{1}{n}}.
	\end{aligned}\right\}
	\!\!\!\implies\!\!\!
	\begin{aligned}
		\frac{1}{r^{2}}
		\diff*{
			\left(
			\frac{r^{2}}{\rho}\diff{P}{r}
			\right)
		}{r}=
		-4\pi G \rho.
		% \text{See the model derivation.}
	\end{aligned}
\end{equation*}\normalsize

Con el objetivo de obtener soluciones numericas para la ecuación Lane-Endem usando métodos eficientes y precisos.

Convertiremos la EDO de segundo orden y de dimensión 1 a una EDO de primer orden y de dimensión 2.
Con ello obtenemos el teorema de Existencia y Unicidad de la ecuación de lane-Endem \cite{diGiovanni2019}, apoyandonos del teorema del punto fijo de Banach.* %podríamos agregar una referencia sobre el teorema del punto fijo de Banach ,quizás Giovanni lo hizo en su paper.

Además utilizaremos el ADM para aproximar la serie analítica \cite{Adomian1995}.Con el ADM obtendremos una serie con polinomios llamados polinomios de Adomiam, que junto el teorema de Taylor* encontramos una cota de error de la solución por el ADM \cite{Umesh2021}.

Nuestra contribución con este trabajo es , aparte de lo ya redactado arriba , es compararla con el método RK5(4)7M ; Utilizando multíples indices politropicos , hecho que verá en las gráficas posteriores.	

Flexible, modular simulation environments are key to many important application fields such as aerospace engineering~\cite{Slotnick2014}, biomedical engineering~\cite{HellersteibBiomed2019}, climate and environmental research~\cite{Climate2020} and many others.
The need to provide smart mathematical and software solutions to combine different aspects of such
simulations in a modular way has been showcased in many publications, e.g.,~\cite{Keyes2013}. With increasing complexity of the respective software environments, the usability and maintainability of the involved software components become a critical issue, which is addressed by a growing research software engineering community (see, e.g., a recent position paper of the German community~\cite{Anzt2021_RSE}).

We present the software package preCICE, which enables black-box coupling of separate solvers for different types of numerical models. It has originally been developed for modular, so-called partitioned, simulations of fluid-structure interactions, i.e., the combination of a flow solver with a structural mechanics solver via a common surface at which forces and displacements are exchanged. Over the past ten years, preCICE has developed into a far more general tool for partitioned simulations, which can handle different types of coupling (weak/strong, explicit/implicit, surface/volume) and any type of equations. Examples range from fluid-structure-acoustics interactions~\cite{Lindner2020_ExaFSA}, over blood flow simulation in the human body~\cite{naseri2020scalable}, free-flow porous media coupling~\cite{Jaust2020}, conjugate heat transfer~\cite{fan2020study}, muscle-tendon system simulations~\cite{Maier2021_Diss}, flow-particle coupling~\cite{besseron2021eulerian}, to coupling between subsurface flow and planning tools for geothermal energy infrastructure~\cite{geoKW}. The coupling is not restricted to a pair of solvers, but has been extended to enable multi-component coupling of arbitrarily many solvers~\cite{Bungartz2014_MIQN_CompuMech}.
preCICE offers comprehensive functionality far beyond simple data exchange: It provides (i) a variety of mapping methods for data transfer between non-matching meshes of different solvers, (ii) quasi-Newton acceleration methods for iterative implicit coupling, and (iii) bottleneck-free point-to-point communication between processes of parallel solvers. preCICE was originally designed with surface coupling in mind, but most features can and have been used for volume coupling as well. All coupling numerics and communication are implemented in a library approach and are fully parallelized. The library can be used via a high-level application programming interface (API) in a minimally-invasive way (from the perspective of the coupled solvers).

The first version of preCICE, presented in~\cite{Gatzhammer2015}, used a server process per coupled solver and was, thus, not very efficient for the coupling of parallel solvers. The authors of~\cite{Uekermann2016, Lindner2019, Bungartz2016_ExaFSA_LNCSE} developed a fully parallel version based on point-to-point communication, which shows good scalability on 10,000s of compute cores. A first overview paper of preCICE was published in 2016~\cite{Bungartz2016_preCICE}, summarizing basic functionality, the API, and the user-specific configuration as well as showing example applications and validation cases. Semantic versioning of preCICE was introduced in 2017.

In this paper, we summarize new developments from 2016 to 2021, i.e., from~\cite{Bungartz2016_preCICE} to the release v2.2, which is part of the first preCICE distribution v2104.0~\cite{Chourdakis2021_Distribution} -- a complete ecosystem of preCICE components. We present an overview of the functionality of preCICE in \autoref{sec:library}, in particular the numerical coupling methods, i.e., quasi-Newton iterations and data mapping. We complement the description of methods for data mapping with a performance and accuracy study using realistic 3D turbine blade meshes. Section~\ref{sec:library}, in addition, gives details on the installation process of preCICE. Beyond the core of preCICE, we describe the newly-developed, ready-to-use adapters~\cite{Uekermann2017_Adapters} for many widely-used simulation software projects in \autoref{sec:adapters}. Section~\ref{sec:tutorials} introduces a simple conjugate heat transfer (CHT) scenario and a simple fluid-structure interaction (FSI) scenario as illustrative examples on how to use preCICE with any of these simulation software projects, followed by the presentation of a systematic multi-level testing infrastructure in \autoref{sec:tests}. In the past years, preCICE has become a widely used software ecosystem, for which we have built up a community of users, as we show in \autoref{sec:community}.

Our description focuses on (i) usability (by providing robust numerical choices, the multitude of ready-to-use adapter codes, well-structured documentation, and easily-accessible illustrative examples), (ii) reliability (by the systematic multi-level testing concept), and (iii) sustainability (by continuous integration, well-defined development and release cycles and a concept to involve the community in the software development). For a more classical description of preCICE including classical validation with benchmarks, we refer the reader to~\cite{Bungartz2016_preCICE}. For an analysis of recent performance and scalability improvements, we refer to~\cite{Bungartz2016_ExaFSA_LNCSE, AminPerf2021}. The contributions of this paper enable new scientific insights in the research fields of our users, but also provide new experiences in scientific software engineering. Results for various applications run with preCICE have been published in many other papers such as~\cite{naseri2020scalable, Jaust2020, fan2020study, andrun2020simulating, Cocco2020} (see also \autoref{sec:community}).

Naturally, preCICE is not the only general-purpose coupling software, which has been developed during the past decades. In the following short summary of related tools, we focus particularly on user-focused and open-source software (i.e., we do not focus on in-house or commercial coupling software, e.g., MpCCI~\cite{Wolf2017}).
There is a number of more multi-scale oriented tools, such as Amuse~\cite{Amuse}, MuMMI~\cite{MuMMI}, MUSCLE 3~\cite{veen2020easing}, MaMiCo~\cite{neumann2016mamico}, or MUI~\cite{tang2015multiscale}. Often, the categories of use cases are not strict. MUI, for example, has recently also been used for fluid-structure interaction~\cite{liu2021parallel}. At the same time, current work on preCICE aims towards certain multi-scale coupling patterns (cf.~\autoref{sec:conclusion}). A good review on multi-scale coupling software is provided in~\cite{groen2019mastering}.
For climate simulations, a number of specialized tools are available, for example OASIS3-MCT~\cite{Craig2017_OASIS3_MCT}, YAC~\cite{Hanke2016_YAC}, and C-Coupler2~\cite{Liu2018_C-Coupler2}.
In principle, the term \textit{coupling} is not well-defined. For example, software such as pyiron~\cite{Janssen2019_Pyiron} or the Kepler Project~\cite{Ludascher2006_Kepler} are referred to as \textit{coupling software} as well, whereas they refer to workflow coupling and not a strong coupling between different simultaneously running simulations.

The two software projects most similar to preCICE are presumably DTK~\cite{Slattery2013_DTK} and OpenPALM~\cite{Duchaine2015_OpenPALM}.
DTK offers an API that targets lower-level operations compared to preCICE. Its main job is to map and communicate data between different meshes in parallel. The implementation of the actual coupling logic is left to the user, which leads to greater flexibility, but also to more development effort for the user. OpenPALM employs a similar API, provided by CWIPI. The \textit{front end} of OpenPALM, however, provides a higher-level API that includes coupling logic. On top, a graphical user interface is available to configure and steer coupled simulations.
The largest difference of preCICE compared to DTK and OpenPALM is presumably the large number of ready-to-use adapters to widely-used simulation software packages (e.g., OpenFOAM, SU2, FEniCS, deal.II, or CalculiX) (cf.~\autoref{sec:adapters}).
